


\section{The kingdom of heaven is like a householder who went out to hire laborers for his vineyard.}
(E. B.)

“Folkebladet,” April 05, 1916

If we attentively try to study the parables or images which Jesus uses in order to present the kingdom of heaven, the kingdom of God, we shall become aware that he very often uses such parables in which this kingdom is presented as work, as activity. And when we look upon his own life here on earth, we cannot avoid becoming aware that he was continually at work, and he was so glad in his work that he even calls it his “food to do the Father’s will.”

It is true that he promises rest to those who come to him, but this never means, as many seem to think, an existence consisting in feeling good, having things easy, and being inactive. No, the rest he promises his own is rest from the unrest of heart and conscience on account of sin. It is the rest which consists in being able, with childlike confidence and blessed joy, to cry: Abba, Father. That rest Jesus gives. But he gives it in order that through it we may receive strength, will, desire, and an inward urge to do something for him.

Jesus himself was no lazy Christian. “He went about doing good,” and the last word he spoke to those who had attached themselves to him was this: “Go into all the world and preach the gospel to all creation.”

In the preaching of God’s word which has now sounded among us for nearly a generation, most has been said about “sin and grace.” And however right and proper and necessary such preaching is, there is nevertheless a deficiency in it. This we can see so clearly from Jesus’ preaching. He too spoke about sin and grace, but he also spoke much—no, most—about works. Look for yourself, and you will be convinced. The reason that our preaching has dwelt almost exclusively upon this great theme, “sin and grace,” is closely grounded in two things. On the one hand, there is the distance we have taken from, and the fear we have acquired of, “dead orthodoxy,” “pure doctrine” without asking after life; and on the other hand, the awakening, partly that by which many of us were more or less seized when we came from Norway, and partly that which passed like a thaw over our people and our congregations some thirteen to twenty years ago. And perhaps also—and not least—the position into which we have come among the Norwegian Lutheran church people here in this country.

Our program has indeed been, and is: the quickening and liberation of the congregations. And quickening and liberation the Lord alone can give to the individual and to his congregation. Therefore it is so natural that the preaching among us has sounded approximately as it is set forth in the programmatic explanation “Seek the Lord,” which by some has been so mercilessly criticized.

For my own part, I have reservations about passing judgment as to whether the criticism mentioned was justified, but I am inclined to hold the opinion that the criticism in question was, at least in part, badly applied. But even if we are clear that the kind of preaching here mentioned is both justified and necessary, and that it is moreover, so to speak, a matter of course among us who have all been more or less influenced by the religious outlook and understanding of the congregation held by the late Prof. Georg Sverdrup, this does not exclude the possibility that there may be one or more deficiencies in it. And I am inclined to believe that perhaps in our preaching in the Free Church we have laid too little weight upon encouraging one another to “love and good works.” It is therefore necessary that we return to this matter again and again, lest we forget it.

Our course as believing individuals and congregations is a course into work—not into self-chosen rest in inactivity.

And that Jesus attaches the greatest importance to this matter we see clearly from his portrayal of the last great day of reckoning, when the King shall determine who shall enter the kingdom and who shall be sent away. Then it is precisely what has been neglected or what has been done that has significance and is decisive.

The work particularly in view here is the more directly Christian work, the work for the salvation of souls and the building up of God’s congregation. And with these expressions our thoughts are at once directed toward what is called mission, foreign mission and home or inner mission.

Some will certainly object that not everyone can take part in that work, that not everyone can go out as missionaries abroad or at home. Yes, that is true. Not everyone can go out, and by no means has everyone received a call to do so. But nevertheless everyone can take part.

In nearly every place among our congregations, work for mission is carried on partly by the congregation as such, and partly by societies within it: mission societies, women’s societies, young people’s societies, children’s societies, and so forth. And anyone whatever can join these societies. Moreover, subscription lists are, so to speak, continually sent out among us, in which contributions are requested for one or another of our common undertakings.

None of us therefore needs to excuse himself by saying that we have no opportunity to do anything. The opportunity we have, but whether we have the will is another matter. Yet there is One of whom it is written that he works in us both to will and to do according to his good pleasure. And if only he is allowed to rule in our heart and our life, we shall become warmer and more zealous friends of our home mission and foreign mission as the days pass and as we consider that many of us, at least, have come far on in the day of our life.

It will become with us as with Paul: for me to live is Christ. He becomes the One whom you desire to please and to whom you desire to be acceptable, the One in whose service you count it your honor and find your joy.

And this inward impulse toward work is the work of the Spirit of God in you. There is certainly much that will hinder us in this work. But he who has appointed us to his service has also given us his gifts in order that we may be made complete for the work of ministry.

In this way it is that the congregation, which is the body of Christ, advances more and more toward the form it is appointed to have, and becomes a means of salvation for the world in which it is set as the light upon the lampstand and the city upon the hill.

Can this goal be reached? Yes, but not until it becomes clear to all true members of the congregation that their calling is work, the work of service, and until they are willing to follow the call.

Here lies the difficulty also in the question of solving the common task we have, which is called our home mission. It will never be solved in the way it ought to be so long as our congregations and individual friends push the thought of it and the responsibility for it away from themselves and seek to quiet their conscience with the consolation that, when they work for foreign mission and the school for pastors, others can take care of this petty business called home mission. No, the one ought to be done, and the other not neglected.
